\input{preamble}

% OK, commencer ici.
%
\begin{document}

\title{Desiderata}

\maketitle

\phantomsection
\label{section-phantom}


\tableofcontents

\section{Introduction}
\label{section-introduction}

\noindent
Il s'agit essentiellement d'une liste de sujets que nous souhaitons ajouter au
projet Champs. Comme nous enrichissons continuellement le projet Champs, cette liste reste
toujours quelque peu en retard sur l'état actuel du projet Champs. En fait, c'était peut-être
une erreur de tenter d'énumérer les sujets à ajouter, car il semble
impossible de la tenir à jour.

\medskip\noindent
Dernière mise à jour : jeudi 31 août 2017.


\section{Conventions}
\label{section-conventions}

\noindent
Nous devrions consacrer un chapitre à une courte liste des conventions employées dans ce document.
Ce chapitre existe déjà, voir
Conventions, section \ref{conventions-section-comments},
mais on pourrait y ajouter bien davantage. Il serait particulièrement utile de repérer
les conventions ``cachées'' et les hypothèses tacites, puis de les y consigner.


\section{Sites et topos}
\label{section-sites}

\noindent
Nous disposons d'un chapitre consacré aux sites et aux faisceaux, voir
Sites, section \ref{sites-section-introduction}.
Nous disposons d'un chapitre consacré aux sites annelés (et aux topos) et aux modules sur ceux-ci, voir
Modules sur les sites, section \ref{sites-modules-section-introduction}.
Nous disposons d'un chapitre consacré à la cohomologie dans ce cadre, voir
Cohomologie sur les sites, section \ref{sites-cohomology-section-introduction}.
Mais on pourrait y ajouter bien davantage, surtout dans le chapitre consacré à la cohomologie.


\section{Champs}
\label{section-stacks}

\noindent
Nous disposons d'un chapitre consacré aux champs (abstraits), voir
Champs, section \ref{stacks-section-introduction}.
Il serait souhaitable :
\begin{enumerate}
\item d'améliorer la présentation de la ``champification'',
\item de donner des exemples de champification,
\item de donner davantage d'exemples en général,
\item d'améliorer la présentation des gerbes.
\end{enumerate}
Voici un exemple de résultat qui n'a pas encore été ajouté : étant donné un faisceau de groupes
abéliens $\mathcal{F}$
sur $\mathcal{C}$, l'ensemble des classes d'équivalence de gerbes de lien
$\mathcal{F}$ est en bijection avec $H^2(\mathcal{C}, \mathcal{F})$.


\section{Méthodes simpliciales}
\label{section-simplicial}

\noindent
Nous disposons d'un chapitre consacré aux méthodes simpliciales, voir
Objets simpliciaux, section \ref{simplicial-section-introduction}.
Il doit être revu et amélioré. La présentation des
rapports entre l'homotopie simpliciale (également appelée
homotopie combinatoire) et les complexes de Kan devrait être améliorée.
Il existe un chapitre consacré aux espaces simpliciaux, voir
Espaces simpliciaux, section \ref{spaces-simplicial-section-introduction}.
Ce chapitre traite brièvement des
espaces topologiques simpliciaux, des sites simpliciaux et des topos simpliciaux.
Nous pouvons développer davantage la ``géométrie algébrique simpliciale'' afin d'étudier
les schémas simpliciaux (ou les espaces algébriques simpliciaux, ou
les champs algébriques simpliciaux) et les questions géométriques, leur cohomologie,
etc.


\section{Cohomologie des schémas}
\label{section-schemes-cohomology}

\noindent
Il existe déjà un chapitre consacré à la cohomologie des faisceaux quasi-cohérents, voir
Cohomologie des schémas, section \ref{coherent-section-introduction}.
Nous disposons d'un chapitre qui traite de la catégorie dérivée des
faisceaux quasi-cohérents sur un schéma, voir
Catégories dérivées des schémas, section \ref{perfect-section-introduction}.
Nous disposons d'un chapitre consacré à la dualité pour les schémas noethériens
et à la dualité relative pour les morphismes de schémas, voir
Dualité pour les schémas, section \ref{duality-section-introduction}.
Nous avons également des chapitres sur la cohomologie \'etale des schémas et sur la
cohomologie cristalline des schémas. Mais la majeure partie du contenu de ces
chapitres est très élémentaire, et on pourrait/devrait y ajouter beaucoup.


\section{Théorie des déformations \`a la Schlessinger}
\label{section-deformation-schlessinger}

\noindent
Nous disposons d'un chapitre consacré à cette matière, voir
Théorie formelle des déformations, section \ref{formal-defos-section-introduction}.
Nous disposons d'un chapitre qui présente des exemples de la théorie générale, voir
Problèmes de déformation, section \ref{examples-defos-section-introduction}.
Nous disposons d'un chapitre, voir
Théorie des déformations, section \ref{defos-section-introduction},
qui traite des déformations d'anneaux (et de modules),
des déformations d'espaces annelés (et de faisceaux de modules),
des déformations de topos annelés (et de faisceaux de modules).
Dans ce chapitre, nous employons le complexe cotangent naïf
pour décrire les obstructions, les déformations du premier ordre et
les automorphismes infinitésimaux. Cette matière a trouvé quelques
applications à l'algébricité des champs de modules dans des chapitres ultérieurs.
Il existe également un chapitre consacré au complexe cotangent complet, voir
Complexe cotangent, section \ref{cotangent-section-introduction}.


\section{Définition des champs algébriques}
\label{section-definition-algebraic-stacks}

\noindent
Un champ algébrique est un champ en groupoïdes sur la catégorie des schémas
munie de la topologie fppf, dont la diagonale est représentable par des
espaces algébriques et qui est la cible d'un morphisme lisse surjectif provenant d'un schéma.
Voir Champs algébriques, section \ref{algebraic-section-algebraic-stacks}.
Un ``champ de Deligne--Mumford'' est un champ algébrique pour lequel il existe
un schéma et un morphisme \'etale surjectif de ce schéma vers le champ,
comme dans l'article \cite{DM} de Deligne et Mumford, voir
Champs algébriques, Définition \ref{algebraic-definition-deligne-mumford}.
Nous réserverons l'expression ``champ d'Artin'' aux champs considérés dans les articles
d'Artin, voir \cite{ArtinI}, \cite{ArtinII} et \cite{ArtinVersal}.
Une définition possible consiste à demander qu'un champ d'Artin soit un champ algébrique
$\mathcal{X}$ au-dessus d'un schéma localement noethérien $S$ tel que
$\mathcal{X} \to S$ soit
localement de type fini\footnote{Ce sont précisément les champs algébriques
sur $S$ qui satisfont les axiomes [-1], [0], [1], [2], [3], [4], [5]
des axiomes d'Artin, section \ref{artin-section-axioms}.}.


\section{Exemples de schémas, d'espaces algébriques et de champs algébriques}
\label{section-examples-stacks}

\noindent
Le projet Champs contient actuellement deux chapitres consacrés aux
champs de modules et à leurs propriétés, voir
Champs de modules, section \ref{moduli-section-introduction} et
Modules des courbes, section \ref{moduli-curves-section-introduction}.
À terme, nous comptons en ajouter d'autres, par exemple :
\begin{enumerate}
\item $\mathcal{A}_g$,
c'est-à-dire les schémas abéliens principalement polarisés de genre $g$,
\item $\mathcal{A}_1 = \mathcal{M}_{1, 1}$, c'est-à-dire
les courbes projectives lisses de genre $1$ à $1$ point marqué,
\item $\mathcal{M}_{g, n}$, c'est-à-dire les courbes projectives lisses de genre $g$
munies de $n$ points marqués deux à deux distincts,
\item $\overline{\mathcal{M}}_{g, n}$, c'est-à-dire les
courbes projectives nodales stables à $n$ points marqués, de genre $g$,
\item $\SheafHom_S(\mathcal{X}, \mathcal{Y})$, le champ de modules des morphismes
(sous des conditions appropriées sur les champs $\mathcal{X}$, $\mathcal{Y}$
et le schéma de base $S$),
\item $\textit{Bun}_G(X) = \SheafHom_S(X, BG)$, le champ des torseurs sous $G$
du programme de Langlands géométrique (sous des conditions appropriées sur le schéma
$X$, le schéma en groupes $G$ et le schéma de base $S$),
\item $\Picardstack_{\mathcal{X}/S}$, c'est-à-dire le champ de Picard associé
à un champ algébrique sur un schéma (ou espace) de base.
\end{enumerate}
Plus généralement, le projet Champs manque quelque peu
d'exemples significatifs du point de vue géométrique.


\section{Propriétés des champs algébriques}
\label{section-stacks-properties}

\noindent
C'est peut-être l'un des projets les plus faciles à entreprendre, puisque l'essentiel de la
théorie élémentaire est maintenant en place. Bien entendu, il s'agit en réalité de propriétés
de morphismes de champs. Nous pouvons définir des singularités (à des facteurs lisses près),
etc. Démontrer qu'un champ normal connexe est irréductible, etc.


\section{Site lisse-\'etale d'un champ algébrique}
\label{section-lisse-etale}

\noindent
Il a été introduit dans
Cohomologie des champs, section \ref{stacks-cohomology-section-lisse-etale}.
Un exemple montrant qu'il n'est pas fonctoriel par rapport aux $1$-morphismes
de champs algébriques est étudié dans
Exemples, section \ref{examples-section-lisse-etale-not-functorial}.
On pourrait bien sûr en dire beaucoup plus, mais il s'avère
très utile de démontrer les résultats à l'aide du ``grand'' site \'etale
dans toute la mesure du possible.



\section{Ce que vous avez toujours voulu savoir sans jamais oser le demander}
\label{section-stacks-fun-lemmas}

\noindent
Il y aura beaucoup de lemmes utiles que l'on emploie sans cesse,
mais qui ne sont pas vraiment énoncés explicitement dans la littérature,
ou pour lesquels il n'est pas aisé de trouver des références. Une boîte à outils.

\medskip\noindent
Exemple : étant donnés deux groupoïdes en schémas $R\Rightarrow U$ et
$R' \Rightarrow U'$, comment exprimer un $1$-morphisme
$[U/R] \to [U'/R']$ uniquement en termes de groupoïdes en schémas ?



\section{Faisceaux quasi-cohérents sur les champs}
\label{section-quasi-coherent}

\noindent
Ils sont définis et étudiés dans le chapitre
Cohomologie des champs, section \ref{stacks-cohomology-section-introduction}.
Les catégories dérivées de modules sont étudiées dans le chapitre
Catégories dérivées des champs, section \ref{stacks-perfect-section-introduction}.
On pourrait enrichir considérablement ces chapitres.



\section{Plat et lisse}
\label{section-flat-smooth}

\noindent
Le théorème d'Artin, selon lequel l'existence d'un morphisme plat surjectif dont la source est un schéma
peut se substituer à la condition d'existence d'un morphisme lisse surjectif, est désormais établi dans
Critères de représentabilité, Théorème \ref{criteria-theorem-bootstrap}.


\section{Théorème de représentabilité d'Artin}
\label{section-representability}

\noindent
Ce sujet est traité dans le chapitre
Axiomes d'Artin, section \ref{artin-section-introduction}.
Nous disposons également d'une application, voir
Quot, Théorème \ref{quot-theorem-coherent-algebraic}.
Il devrait y avoir bien davantage d'applications, et le chapitre
lui-même doit également être remanié.


\section{Les champs DM admettent des morphismes finis surjectifs provenant de schémas}
\label{section-dm-finite-cover}

\noindent
Nous disposons déjà du résultat correspondant pour les espaces algébriques, voir
Limites d'espaces, section \ref{spaces-limits-section-finite-cover}.
Il manque le résultat pour les champs DM et quasi-DM.


\section{Article de Martin Olsson sur la propreté}
\label{section-proper-parametrization}

\noindent
Cet article démontre que deux notions de propreté coïncident. La première partie de ce résultat
est désormais disponible sous la forme du lemme de Chow pour les champs algébriques, voir
Compléments sur les morphismes de champs, Théorème
\ref{stacks-more-morphisms-theorem-chow-finite-type}.
Il en résulte que, dans certains cas, il suffit de recourir aux anneaux de valuation discrète
pour vérifier le critère valuatif de propreté
des champs algébriques; voir
Compléments sur les morphismes de champs, section
\ref{stacks-more-morphisms-section-Noetherian-valuative-criterion}.


\section{Image directe de faisceaux cohérents par un morphisme propre}
\label{section-proper-pushforward}

\noindent
Nous pouvons maintenant commencer à étudier cette question, puisque nous disposons du lemme de Chow pour
les champs algébriques; voir la section précédente.


\section{Keel et Mori}
\label{section-keel-mori}

\noindent
Voir \cite{K-M}. Leur résultat a été ajouté dans
Compléments sur les morphismes de champs, section
\ref{stacks-more-morphisms-section-Keel-Mori}.


\section{Ajouter d'autres sujets ici}
\label{section-add-more}

\noindent
En réalité, non : nous n'aurions jamais dû commencer cette liste dans le
projet Champs lui-même ! Il existe ailleurs une liste de tâches
qui est bien plus facile à mettre à jour.


\begin{multicols}{2}[\section{Autres chapitres}]
\noindent
Préliminaires
\begin{enumerate}
\item \hyperref[introduction-section-phantom]{Introduction}
\item \hyperref[conventions-section-phantom]{Conventions}
\item \hyperref[sets-section-phantom]{Théorie des ensembles}
\item \hyperref[categories-section-phantom]{Catégories}
\item \hyperref[topology-section-phantom]{Topologie}
\item \hyperref[sheaves-section-phantom]{Faisceaux sur les espaces}
\item \hyperref[sites-section-phantom]{Sites et faisceaux}
\item \hyperref[stacks-section-phantom]{Champs}
\item \hyperref[fields-section-phantom]{Corps}
\item \hyperref[algebra-section-phantom]{Algèbre commutative}
\item \hyperref[brauer-section-phantom]{Groupes de Brauer}
\item \hyperref[homology-section-phantom]{Algèbre homologique}
\item \hyperref[derived-section-phantom]{Catégories dérivées}
\item \hyperref[simplicial-section-phantom]{Méthodes simpliciales}
\item \hyperref[more-algebra-section-phantom]{Compléments d'algèbre}
\item \hyperref[smoothing-section-phantom]{Lissage des morphismes d'anneaux}
\item \hyperref[modules-section-phantom]{Faisceaux de modules}
\item \hyperref[sites-modules-section-phantom]{Modules sur les sites}
\item \hyperref[injectives-section-phantom]{Objets injectifs}
\item \hyperref[cohomology-section-phantom]{Cohomologie des faisceaux}
\item \hyperref[sites-cohomology-section-phantom]{Cohomologie sur les sites}
\item \hyperref[dga-section-phantom]{Algèbre différentielle graduée}
\item \hyperref[dpa-section-phantom]{Algèbre à puissances divisées}
\item \hyperref[sdga-section-phantom]{Faisceaux différentiels gradués}
\item \hyperref[hypercovering-section-phantom]{Hyperrecouvrements}
\end{enumerate}
Schémas
\begin{enumerate}
\setcounter{enumi}{25}
\item \hyperref[schemes-section-phantom]{Schémas}
\item \hyperref[constructions-section-phantom]{Constructions de schémas}
\item \hyperref[properties-section-phantom]{Propriétés des schémas}
\item \hyperref[morphisms-section-phantom]{Morphismes de schémas}
\item \hyperref[coherent-section-phantom]{Cohomologie des schémas}
\item \hyperref[divisors-section-phantom]{Diviseurs}
\item \hyperref[limits-section-phantom]{Limites de schémas}
\item \hyperref[varieties-section-phantom]{Variétés}
\item \hyperref[topologies-section-phantom]{Topologies sur les schémas}
\item \hyperref[descent-section-phantom]{Descente}
\item \hyperref[perfect-section-phantom]{Catégories dérivées de schémas}
\item \hyperref[more-morphisms-section-phantom]{Compléments sur les morphismes}
\item \hyperref[flat-section-phantom]{Compléments sur la platitude}
\item \hyperref[groupoids-section-phantom]{Schémas en groupoïdes}
\item \hyperref[more-groupoids-section-phantom]{Compléments sur les schémas en groupoïdes}
\item \hyperref[etale-section-phantom]{Morphismes étales de schémas}
\end{enumerate}
Thèmes de théorie des schémas
\begin{enumerate}
\setcounter{enumi}{41}
\item \hyperref[chow-section-phantom]{Homologie de Chow}
\item \hyperref[intersection-section-phantom]{Théorie de l'intersection}
\item \hyperref[pic-section-phantom]{Schémas de Picard des courbes}
\item \hyperref[weil-section-phantom]{Théories cohomologiques de Weil}
\item \hyperref[adequate-section-phantom]{Modules adéquats}
\item \hyperref[dualizing-section-phantom]{Complexes dualisants}
\item \hyperref[duality-section-phantom]{Dualité pour les schémas}
\item \hyperref[discriminant-section-phantom]{Discriminants et différentes}
\item \hyperref[derham-section-phantom]{Cohomologie de de Rham}
\item \hyperref[local-cohomology-section-phantom]{Cohomologie locale}
\item \hyperref[algebraization-section-phantom]{Géométrie algébrique et formelle}
\item \hyperref[curves-section-phantom]{Courbes algébriques}
\item \hyperref[resolve-section-phantom]{Résolution des surfaces}
\item \hyperref[models-section-phantom]{Réduction semi-stable}
\item \hyperref[functors-section-phantom]{Foncteurs et morphismes}
\item \hyperref[equiv-section-phantom]{Catégories dérivées de variétés}
\item \hyperref[pione-section-phantom]{Groupes fondamentaux de schémas}
\item \hyperref[etale-cohomology-section-phantom]{Cohomologie étale}
\item \hyperref[crystalline-section-phantom]{Cohomologie cristalline}
\item \hyperref[proetale-section-phantom]{Cohomologie pro-étale}
\item \hyperref[relative-cycles-section-phantom]{Cycles relatifs}
\item \hyperref[more-etale-section-phantom]{Compléments de cohomologie étale}
\item \hyperref[trace-section-phantom]{La formule des traces}
\end{enumerate}
Espaces algébriques
\begin{enumerate}
\setcounter{enumi}{64}
\item \hyperref[spaces-section-phantom]{Espaces algébriques}
\item \hyperref[spaces-properties-section-phantom]{Propriétés des espaces algébriques}
\item \hyperref[spaces-morphisms-section-phantom]{Morphismes d'espaces algébriques}
\item \hyperref[decent-spaces-section-phantom]{Espaces algébriques décents}
\item \hyperref[spaces-cohomology-section-phantom]{Cohomologie des espaces algébriques}
\item \hyperref[spaces-limits-section-phantom]{Limites d'espaces algébriques}
\item \hyperref[spaces-divisors-section-phantom]{Diviseurs sur les espaces algébriques}
\item \hyperref[spaces-over-fields-section-phantom]{Espaces algébriques sur des corps}
\item \hyperref[spaces-topologies-section-phantom]{Topologies sur les espaces algébriques}
\item \hyperref[spaces-descent-section-phantom]{Descente et espaces algébriques}
\item \hyperref[spaces-perfect-section-phantom]{Catégories dérivées d'espaces}
\item \hyperref[spaces-more-morphisms-section-phantom]{Compléments sur les morphismes d'espaces}
\item \hyperref[spaces-flat-section-phantom]{Platitude sur les espaces algébriques}
\item \hyperref[spaces-groupoids-section-phantom]{Groupoïdes dans les espaces algébriques}
\item \hyperref[spaces-more-groupoids-section-phantom]{Compléments sur les groupoïdes dans les espaces}
\item \hyperref[bootstrap-section-phantom]{Amorçage}
\item \hyperref[spaces-pushouts-section-phantom]{Sommes amalgamées d'espaces algébriques}
\end{enumerate}
Thèmes de géométrie
\begin{enumerate}
\setcounter{enumi}{81}
\item \hyperref[spaces-chow-section-phantom]{Groupes de Chow des espaces}
\item \hyperref[groupoids-quotients-section-phantom]{Quotients de groupoïdes}
\item \hyperref[spaces-more-cohomology-section-phantom]{Compléments sur la cohomologie des espaces}
\item \hyperref[spaces-simplicial-section-phantom]{Espaces simpliciaux}
\item \hyperref[spaces-duality-section-phantom]{Dualité pour les espaces}
\item \hyperref[formal-spaces-section-phantom]{Espaces algébriques formels}
\item \hyperref[restricted-section-phantom]{Algébrisation des espaces formels}
\item \hyperref[spaces-resolve-section-phantom]{Résolution des surfaces revisitée}
\end{enumerate}
Théorie des déformations
\begin{enumerate}
\setcounter{enumi}{89}
\item \hyperref[formal-defos-section-phantom]{Théorie formelle des déformations}
\item \hyperref[defos-section-phantom]{Théorie des déformations}
\item \hyperref[cotangent-section-phantom]{Le complexe cotangent}
\item \hyperref[examples-defos-section-phantom]{Problèmes de déformation}
\end{enumerate}
Champs algébriques
\begin{enumerate}
\setcounter{enumi}{93}
\item \hyperref[algebraic-section-phantom]{Champs algébriques}
\item \hyperref[examples-stacks-section-phantom]{Exemples de champs}
\item \hyperref[stacks-sheaves-section-phantom]{Faisceaux sur les champs algébriques}
\item \hyperref[criteria-section-phantom]{Critères de représentabilité}
\item \hyperref[artin-section-phantom]{Axiomes d'Artin}
\item \hyperref[quot-section-phantom]{Espaces de Quot et de Hilbert}
\item \hyperref[stacks-properties-section-phantom]{Propriétés des champs algébriques}
\item \hyperref[stacks-morphisms-section-phantom]{Morphismes de champs algébriques}
\item \hyperref[stacks-limits-section-phantom]{Limites de champs algébriques}
\item \hyperref[stacks-cohomology-section-phantom]{Cohomologie des champs algébriques}
\item \hyperref[stacks-perfect-section-phantom]{Catégories dérivées de champs}
\item \hyperref[stacks-introduction-section-phantom]{Introduction aux champs algébriques}
\item \hyperref[stacks-more-morphisms-section-phantom]{Compléments sur les morphismes de champs}
\item \hyperref[stacks-geometry-section-phantom]{La géométrie des champs}
\end{enumerate}
Thèmes de théorie des espaces de modules
\begin{enumerate}
\setcounter{enumi}{107}
\item \hyperref[moduli-section-phantom]{Champs de modules}
\item \hyperref[moduli-curves-section-phantom]{Espaces de modules de courbes}
\end{enumerate}
Divers
\begin{enumerate}
\setcounter{enumi}{109}
\item \hyperref[examples-section-phantom]{Exemples}
\item \hyperref[exercises-section-phantom]{Exercices}
\item \hyperref[guide-section-phantom]{Guide bibliographique}
\item \hyperref[desirables-section-phantom]{Desiderata}
\item \hyperref[coding-section-phantom]{Style de programmation}
\item \hyperref[obsolete-section-phantom]{Obsolète}
\item \hyperref[fdl-section-phantom]{Licence de documentation libre GNU}
\item \hyperref[index-section-phantom]{Index généré automatiquement}
\end{enumerate}
\end{multicols}



\bibliography{my}
\bibliographystyle{amsalpha}

\end{document}
